% C084 | 双曲线焦点三角形的内角平分线（简要源稿；无 TikZ 图）
% 本文件是知识点简介片段，不是可独立编译的 LaTeX 主文档。
% 详见 01_statement--06_summary；排版由 main.tex 接入项目模板。
\textbf{C084\quad 双曲线焦点三角形的内角平分线}

设双曲线
\[
H:\frac{x^2}{a^2}-\frac{y^2}{b^2}=1\quad(a,b>0),
\qquad c^2=a^2+b^2,
\]
左右焦点 $F_1(-c,0)$、$F_2(c,0)$。
取双曲线上非顶点 $P(x_0,y_0)$（$y_0\ne0$），
$\angle F_1PF_2$ 的\textbf{内角平分线}交 $x$ 轴于 $M(m,0)$。
则
\[
\boxed{x_0m=a^2,\qquad M\left(\frac{a^2}{x_0},0\right)}.
\]
更进一步，$PM$ 正是双曲线在 $P$ 处的\textbf{切线}，其方程为
\[
\boxed{\frac{x_0x}{a^2}-\frac{y_0y}{b^2}=1}.
\]
\textbf{最低门槛思路：}角平分线定理给出
$m=c(PF_1-PF_2)/(PF_1+PF_2)$，
距离平方差给出 $PF_1^2-PF_2^2=4cx_0$，
再用双曲线定义 $|PF_1-PF_2|=2a$ 即得结论。

\textbf{几何补充：}焦点三角形的内切圆与焦点轴的切点为
$T=(a,0)$（右支）或 $T=(-a,0)$（左支），
它一般\textbf{不是}角平分线的轴上交点 $M$。
